\documentclass[10pt]{article}
\usepackage[T1]{fontenc}
\usepackage[utf8]{inputenc}
\usepackage[french]{babel}
\usepackage{amsmath,amssymb,mathrsfs,mathpazo}
\usepackage{geometry}
\geometry{a4paper,margin=2cm}
\usepackage{graphicx,array,tabularx,tabularray,colortbl,multirow,multicol,booktabs,enumitem,pifont,ulem,fancybox,diagbox,cancel}
\usepackage{pstricks,pstricks-add,pst-plot,pst-tree,pst-node,pst-text,pst-func,pst-math,pst-bspline,pst-3dplot,pst-eucl,pst-solides3d}
\let\pscalculate\undefined
\usepackage{pst-bezier}
\usepackage{tikz,pgfplots,tkz-tab}
\usetikzlibrary{arrows,arrows.meta,patterns,calc,positioning,trees,decorations.markings,shapes.geometric,shapes.misc,shapes.arrows,intersections,decorations.pathmorphing,decorations.pathreplacing,backgrounds,shadows,babel}
\pgfplotsset{compat=1.18}
\usepgfplotslibrary{fillbetween}
\usepackage[np]{numprint}
\usepackage{listings,makecell,pas-tableur,twoopt}
\let\barre\undefined
\usepackage[pstricks]{tablvar}
\definecolor{darkgreen}{RGB}{0,100,0}
\providecommand{\red}{\color{red}}
\providecommand{\blue}{\color{blue}}
\providecommand{\green}{\color{green}}
\providecommand{\black}{\color{black}}
\providecommand{\white}{\color{white}}
\providecommand{\gray}{\color{gray}}
\providecommand{\cyan}{\color{cyan}}
\providecommand{\magenta}{\color{magenta}}
\providecommand{\yellow}{\color{yellow}}
\providecommand{\violet}{\color{violet}}
\providecommand{\orange}{\color{orange}}
\providecommand{\euro}{\texteuro}
\providecommand{\EUR}{\texteuro}
\providecommand{\ieme}{\textsuperscript{e}}
\providecommand{\degres}{\ensuremath{^\circ}}
\providecommand{\degree}{\ensuremath{^\circ}}
\providecommand{\vv}[1]{\overrightarrow{#1}}
\providecommand{\Coord}[1]{#1}
\providecommand{\decosix}{}
\providecommand{\starredbullet}{\ensuremath{\bullet}}
\pagestyle{empty}
\setlength{\parindent}{0pt}
\usepackage{adjustbox}\DeclareUnicodeCharacter{2236}{:}\DeclareUnicodeCharacter{2190}{\ensuremath{\leftarrow}}
\providecommand{\R}{\mathbb{R}}
\providecommand{\N}{\mathbb{N}}
\providecommand{\Z}{\mathbb{Z}}
\providecommand{\Q}{\mathbb{Q}}
\providecommand{\D}{\mathbb{D}}
\providecommand{\C}{\mathbb{C}}
\providecommand{\e}{{\mathrm{e}}}
\providecommand{\vect}[1]{\overrightarrow{#1}}
\providecommand{\vectt}[1]{\overrightarrow{\mathrm{#1}}}
\providecommand{\vvt}[1]{\overrightarrow{\mathrm{#1}}}
\providecommand{\np}[2][]{{#2~#1}}
\providecommand{\vv}[1]{{\overrightarrow{#1}}}
\renewcommand{\d}{{\mathrm{d}}}
\providecommand{\degres}{{^\circ}}
\providecommand{\degre}{{^\circ}}
\providecommand{\blue}{\color{blue}}
\providecommand{\red}{\color{red}}
\providecommand{\enskip}{\quad}
\providecommand{\llbracket}{\lbrack\!\lbrack}
\providecommand{\rrbracket}{\rbrack\!\rbrack}
\providecommand{\euro}{{\text{€}}}
\providecommand{\no}{{\text{n°}}}
\providecommand{\interval}[2]{{[#1;#2]}}
\providecommand{\intervFO}[2]{{[#1;#2[}}
\providecommand{\dotfill}{\ldots}

\providecommand{\ds}{\displaystyle}
\providecommand{\textding}[1]{#1}
\providecommand{\Oij}{\(\mathrm{O};\vec{\imath},\vec{\jmath}\)}
\providecommand{\Oijk}{\(\mathrm{O};\vec{\imath},\vec{\jmath},\vec{k}\)}
\providecommand{\Ouv}{\(\mathrm{O};\vec{u},\vec{v}\)}
\providecommand{\intervFF}[2]{\ensuremath{[#1;#2]}}
\renewcommand{\euro}{\ifmmode\text{\texteuro}\else\texteuro\fi}
\DeclareUnicodeCharacter{00B0}{\ifmmode^{\circ}\else\textdegree\fi}
\DeclareUnicodeCharacter{20AC}{\ifmmode\text{\texteuro}\else\texteuro\fi}
% Énoncés : Éducation nationale. Transcription et corrigés : APMEP, Denis Vergès et contributeurs. Découpage et classement : CoopMaths, Sébastien Lozano, Éric Elter et contributeurs. Adaptation de la présentation : Maths Explorer.
% CC BY-SA : https://coopmaths.fr/www/about/
% https://coopmaths.fr/alea/static/e3c/2020/tex/e3c_2020_00_sujet11_2_cor.tex
\begin{document}
\noindent\textit{APMEP / CoopMaths. CC BY-SA. Présentation adaptée par Maths Explorer.}\par\medskip
\subsection*{1. Le triangle $ABC$ est-il isocèle en $B$ ?}
On a $BA^2 = (1 - (-1))^2 + (2 - (-3))^2 = 4 + 25 = 29$ ; $BC^2 = (7 - 1)^2 + (1 - 2)^2 = 36 + 1 = 37$, donc $BA^2 \neq BC^2$ : le triangle n'est pas isocèle en $B$.
	
	\subsection*{2. Déterminer la valeur arrondie au dixième de degré de l'angle $\widehat{BAC}$.}
Par définition du produit scalaire :
	\[
	\overrightarrow{AB} \cdot 	\overrightarrow{AC} = AB \times AC \times \cos(\widehat{BAC})
	\]
	
	Avec $	\overrightarrow{AB}\displaystyle\binom{2}{5}$, $	\overrightarrow{AC}\displaystyle\binom{8}{4}$, $AB = \sqrt{29}$ et $AC = \sqrt{80}$, l'égalité devient :
	\[
	16 + 20 = \sqrt{29} \times \sqrt{80} \times \cos(\widehat{BAC}) \quad \text{d'où} \quad \cos(\widehat{BAC}) = \dfrac{36}{\sqrt{29 \times 80}}
	\]
	
	La calculatrice donne $\widehat{BAC} \approx 41,6^\circ$.
	
	\subsection*{3. On considère le point $H$ de coordonnées $(2,6 ; -1,2)$. Le point $H$ est-il le projeté orthogonal du point $B$ sur la droite $(AC)$ ?}
	
	Le point $H$ est le projeté orthogonal du point $B$ sur la droite $(AC)$ si :
	\begin{itemize}
		\item $H$ appartient à la droite $(AC)$ ;
		\item $(BH)$ est perpendiculaire à $(AC)$.
	\end{itemize}
	
	Or $\iff{AH}(3,6 ; 1,8)$ et $\iff{AC}(8 ; 4)$ : manifestement ces vecteurs sont colinéaires : $H$ appartient à la droite $(AC)$.
	
	De plus $\iff{AC}(8 ; 4)$ et $\iff{BH}(1,6 ; -3,2)$, d'où :
	\[
	\iff{AC} \cdot \iff{BH} = 8 \times 1,8 - 4 \times 3,6 = 14,4 - 14,4 = 0
	\]
	
	Les vecteurs sont orthogonaux, donc le point $H$ est le projeté orthogonal du point $B$ sur la droite $(AC)$.

\end{document}
